\documentclass[12pt]{article}
\usepackage[utf8]{inputenc}
\usepackage[T1]{fontenc}
\usepackage{hyperref}
\title{Waioli Language Resources for K-12 Curriculum Development}
\author{Compiled by Claude Sonnet 5.5 (claude-sonnet-5-5) for Andr\'as Kornai}
\date{}
\begin{document}
\maketitle

\section{Introduction}
Waioli (ISO 639-3 wli) is a North Halmahera language of Maluku Utara, Indonesia. Glottolog groups Sahu and Waioli together in the Sahuan branch of North Halmahera, apart from the Galela--Loloda group. The search found very little Waioli-specific material. It found no Waioli grammar, dictionary, school materials, audio or speaker count.

\section{Primary Educational Resources}
No primers, textbooks or school materials for Waioli were found in search results.

\section{Linguistic and Reference Materials}
\begin{itemize}
\item An Indonesian-language lexicostatistical study in \emph{Jurnal Penelitian Humano} (2021, Universitas Khairun) on the kinship of non-Austronesian languages of West Halmahera, including Waioli, Sahu, Gamkonora, Ibo, Tobaru and Loloda: \url{https://ejournal.unkhair.ac.id/index.php/humano/article/download/3281/2434}
\end{itemize}

\section{Religious and Community Resources}
None were found in search results.

\section{Digital and Audio Resources}
None for Waioli were found in search results. The PARADISEC and ASJP results were for neighbouring languages (Loloda, Galela), not Waioli, and are not listed.

\section{Curriculum Development Context}
Documentation of Waioli appears to be minimal. Any curriculum would first require basic documentation: a wordlist, a description and recordings, made with the community and in cooperation with regional universities such as Universitas Khairun in Ternate.
\end{document}